% chapters/01_signals_as_vectors.tex
%
% Part I, Chapter 1.  Adapted and rewritten from "A Geometric Tour of
% Machine Learning" Chapter 1 (Mathematical Foundation), with a new
% musically-motivated opening, tighter prose, and margin notes to keep the
% reader oriented.

\chapter{Signals as vectors}
\labch{signals_as_vectors}

\epigraph{The fool doth think he is wise, but the wise man knows himself
to be a fool.}{William Shakespeare, \emph{As You Like It}}

%------------------------------------------------------------------------
\section{Two notes}
%------------------------------------------------------------------------

\noindent Listen to a sustained violin note and a sustained flute note,
played at the same pitch.  They have the same fundamental frequency.
They sound nothing alike.  Now play them together.  They blend into
something that is neither violin nor flute --- a third instrument that
exists only while they are sounding together.

\marginnote{%
This is the setup.  The math we develop in this chapter --- vector spaces,
inner products, projections --- is the language in which the difference
between a violin and a flute, and the meaning of ``blend,'' can be made
precise.}

The question \emph{why do they sound different} can be answered at two
levels.  At the musical level the answer is timbre: the way each
instrument's overtones reinforce or attenuate the others, the attack
transient, the bow noise or the breath noise, the small fluctuations in
pitch and amplitude.  At the mathematical level the answer is that the
two recordings, considered as long sequences of numbers, point in
different directions in a very high-dimensional space.  When we add them,
we are doing geometry: the result is a third point in the same space,
neither of the originals but related to both.

This book lives in the gap between those two answers.  The technical
chapters (Part~\textsc{II}) build software that operates on sounds as
high-dimensional vectors.  The creative chapters (Part~\textsc{III})
reason about timbre, style, and composition in those terms.  This first
part is about establishing the geometry --- the structure of the space
that makes the technical work possible and the creative work coherent.

The central claim is simple, and it will reappear in many guises:
\emph{representing a signal means projecting it onto a set of reference
vectors}.  The Fourier transform projects onto sinusoids.  The wavelet
transform projects onto localized oscillations.  Principal component
analysis projects onto data-derived directions of maximum variance.  A
neural network projects onto weight vectors learned from data, then
applies a nonlinearity and projects again.  The mathematics underlying
all four is the same: an inner product, repeated.  What differs is the
choice of reference vectors and whether they are designed by hand or
learned from data.

\marginnote{%
\textbf{Notation.}  Vectors are denoted by bold lowercase letters
$\xb$, $\yb$.  Matrices by bold uppercase, $\Wb$, $\mathbf{B}$.  The
$n$-th component of vector $\xb$ is written $\xb[n]$.  The inner product
of two vectors is $\inner{\xb}{\yb}$, and convolution is $\xb \conv \yb$.
Appendix~\ref{ch:matrix_theory} reviews matrix conventions.}

%------------------------------------------------------------------------
\section{Vector spaces}
%------------------------------------------------------------------------

To think of a signal as a point in space, we need to say precisely what
``space'' means.  A \emph{vector space} is a set $V$ on which two
operations are defined --- addition of elements, and multiplication by
scalars --- which satisfy a list of axioms.\sidenote{The full axiom list
appears in Appendix~\ref{ch:matrix_theory}; for our purposes, what matters
is that they let us add signals and scale them without leaving the space.}
The elements of $V$ are called \emph{vectors}.  When the scalar field is
the real numbers $\R$, we say $V$ is a real vector space; when it is the
complex numbers $\C$, a complex vector space.

A discrete-time signal of length $N$ is naturally an element of $\R^N$
(or $\C^N$, if we allow complex samples).  The $n$-th component is the
$n$-th sample.  Adding two signals samplewise is vector addition.
Multiplying every sample by a constant is scalar multiplication.  The
zero vector is the silent signal.

In this book we work exclusively in finite dimensions.  The infinite-dim\-en\-sion\-al
theory of Hilbert spaces is beautiful and important, but everything we
need can be done in $\R^N$ or $\C^N$, and the finite case is what the
computer manipulates anyway.

%------------------------------------------------------------------------
\section{Norm: the size of a signal}
%------------------------------------------------------------------------

To say one signal is louder than another, or that one is close to
another, we need a notion of size and of distance.  A \emph{norm} on a
vector space $V$ is a function $\norm{\cdot}: V \to \R$ satisfying
\emph{(i)} $\norm{\xb} \ge 0$ with equality only at $\xb = \mathbf{0}$,
\emph{(ii)} $\norm{\alpha \xb} = \abs{\alpha} \norm{\xb}$ for any scalar
$\alpha$, and \emph{(iii)} the triangle inequality $\norm{\xb + \yb} \le
\norm{\xb} + \norm{\yb}$.

The norm we will use almost everywhere is the \emph{Euclidean} or
\emph{$L^2$} norm,
\begin{equation}
    \norm{\xb}_2 \;=\; \sqrt{\sum_{n=0}^{N-1} \abs{\xb[n]}^2}.
    \label{eq:l2norm}
\end{equation}
The square of this quantity, $\norm{\xb}_2^2$, has a physical
interpretation when $\xb$ is an audio signal: it is the \emph{energy} of
the signal.  We will often work with energy directly rather than with
the norm, because it avoids the square root.

\marginnote{%
Other norms exist and matter.  The $L^1$ norm $\sum_n \abs{\xb[n]}$ is
used in sparse coding because it prefers solutions with many zeros.  The
$L^\infty$ norm $\max_n \abs{\xb[n]}$ is the peak amplitude.  We will
meet $L^1$ again in Part~\textsc{III} when we discuss matching pursuit.}

A vector space equipped with a norm is called a \emph{Banach space}.
Banach spaces are what we need to talk about size and convergence.  We
need one more piece of structure --- an inner product --- to talk about
angles and projections.  That promotes a Banach space to a Hilbert space.

%------------------------------------------------------------------------
\section{Inner product: global similarity}
%------------------------------------------------------------------------

The inner product is the single most important operation in this book.
For two real-valued signals $\xb$ and $\yb$ of length $N$,
\begin{equation}
    \inner{\xb}{\yb} \;=\; \sum_{n=0}^{N-1} \xb[n]\, \yb[n].
    \label{eq:inner}
\end{equation}
For complex-valued signals, the second factor is conjugated:
$\inner{\xb}{\yb} = \sum_n \xb[n]\, \overline{\yb[n]}$.

Geometrically, the inner product measures \emph{similarity}.  The clean
statement is the following.  Two vectors with the same magnitude have a
large inner product when they point in the same direction, zero when
they are perpendicular, and a large negative value when they point in
opposite directions.\sidenote{The relation
$\inner{\xb}{\yb} = \norm{\xb}\, \norm{\yb} \cos\theta$, where $\theta$
is the angle between them, makes this precise.  When $\theta = 0$ the
cosine is $1$; when $\theta = \pi/2$ it is zero; when $\theta = \pi$ it
is $-1$.}  When two vectors are identical, their inner product is the
square of their norm: $\inner{\xb}{\xb} = \norm{\xb}^2$.  When they are
orthogonal ($\inner{\xb}{\yb} = 0$), they share nothing along each
other's direction.

The inner product is \emph{global}.  Every sample of $\xb$ is multiplied
with every corresponding sample of $\yb$, and the products are summed
into a single number.  No spatial or temporal locality is preserved ---
the contribution of a sample at the beginning of the signal is
indistinguishable, in the final sum, from the contribution of a sample
at the end.  This is sometimes what we want (when comparing the overall
frequency content of two signals, say) and sometimes not (when looking
for a specific pattern at an unknown time).  When we want locality, we
use convolution instead, which is the next section's subject.

A vector space equipped with an inner product (and the norm derived
from it via $\norm{\xb} = \sqrt{\inner{\xb}{\xb}}$) is called a
\emph{Hilbert space}.  Hilbert spaces are the natural setting for
analysis, synthesis, and learning.  Every space we work with in this book
will be a Hilbert space.

%------------------------------------------------------------------------
\section{Convolution: local similarity}
\label{sec:convolution}
%------------------------------------------------------------------------

A signal is rarely interesting in its entirety.  When we listen, we
attend to local features --- the attack of a note, a particular harmonic
fluctuation, a pause.  The mathematical operation that picks out local
features is the \emph{convolution}.

Given a signal $\xb$ and a short pattern $\mathbf{h}$ (called the
\emph{kernel} or \emph{filter}) of length $M$, the convolution
$\xb \conv \mathbf{h}$ is the sequence
\begin{equation}
    (\xb \conv \mathbf{h})[n]
    \;=\;
    \sum_{m=0}^{M-1} \xb[n+m]\, \mathbf{h}[m].
    \label{eq:conv}
\end{equation}
For each output position $n$, the kernel is aligned with the signal
starting at sample $n$, and the inner product is taken over just those
$M$ samples.  This is the local version of the inner product from the
previous section: instead of computing one number that summarises the
entire signal, we compute many numbers, each summarising a small
neighbourhood.

\marginnote{%
A small example.  If $\mathbf{h} = (1, -1)$, then
$(\xb \conv \mathbf{h})[n] = \xb[n+1] - \xb[n]$, the discrete derivative.
If $\mathbf{h} = (1, 1)/2$, we get a two-point smoother.  Choosing
$\mathbf{h}$ chooses what features the convolution highlights.  This is
precisely what a convolutional neural network learns to do: its kernels
are exactly such patterns, but discovered from data rather than designed.}

Convolution sits at the centre of both signal processing and modern
machine learning.  In signal processing it underlies filtering,
reverberation, pattern detection, and edge detection.  In machine
learning it is the basis of convolutional neural networks, which we
build in Part~\textsc{II}.  In both fields, convolution is doing the
same thing: \emph{repeated inner products, slid along the signal}.

The kernel $\mathbf{h}$ is itself a signal --- a short one, typically.
Choosing the kernel is choosing what local feature to look for.  A
kernel that resembles a quarter-period of a sinusoid will respond
strongly where the signal looks like that fragment of sinusoid; this is
the seed of Fourier analysis.  A kernel that resembles an edge will
respond at sharp transitions; this is the seed of the Sobel filter and
of the first layer of an image conv-net.  Different kernels reveal
different features; the right kernel to use depends on what we are
trying to discover.

There is a deep theorem connecting convolution to Fourier analysis,
namely that convolution in the time domain is equivalent to pointwise
multiplication in the frequency domain.  We will state and prove it in
Chapter~\ref{ch:apriori} when we have set up the Fourier transform;
Appendix~\ref{ch:convolution_theorem} contains a fuller proof.  The
practical implication is that convolution can be computed quickly via
the FFT, which is why convolutional networks scale.

%------------------------------------------------------------------------
\section{Basis and projection}
%------------------------------------------------------------------------

We have inner products (global similarity) and convolutions (local
similarity).  The next step is to use them to \emph{decompose} a signal
into pieces.

Given a set of vectors $\{\mathbf{b}_1, \mathbf{b}_2, \ldots,
\mathbf{b}_K\}$ in a vector space $V$, the \emph{span} of the set is the
collection of all linear combinations
\(
\sum_k \alpha_k \mathbf{b}_k
\)
that can be formed from them.  If every vector in $V$ can be written
this way --- if the span of the set is all of $V$ --- and if no vector
in the set is a linear combination of the others, the set is called a
\emph{basis} for $V$.

A basis is the minimal set of building blocks for the space.  In $\R^N$
there are always exactly $N$ vectors in any basis; we can pick the
``natural'' basis $\mathbf{e}_n$ (the $n$-th unit vector, equal to one
at position $n$ and zero elsewhere) or any other set of $N$ linearly
independent vectors.

The point of having a basis is that any vector $\xb \in V$ can be
written uniquely as
\begin{equation}
    \xb \;=\; \sum_{k=1}^{N} c_k\, \mathbf{b}_k.
    \label{eq:expansion}
\end{equation}
The numbers $c_k$ are the \emph{coordinates} of $\xb$ in the basis
$\{\mathbf{b}_k\}$.  They are what change when we change basis; the
underlying vector $\xb$ does not.

When the basis is \emph{orthonormal} --- meaning the basis vectors are
mutually perpendicular and each has unit norm --- there is a beautifully
simple formula for the coordinates:
\begin{equation}
    c_k \;=\; \inner{\xb}{\mathbf{b}_k}.
    \label{eq:coord_orth}
\end{equation}
The coordinate $c_k$ is just the inner product of $\xb$ with the $k$-th
basis vector.  Combining this with~\eqref{eq:expansion},
\begin{equation}
    \xb \;=\; \sum_{k=1}^{N} \inner{\xb}{\mathbf{b}_k}\, \mathbf{b}_k.
    \label{eq:reconstruction}
\end{equation}
This is the single most important formula in classical signal processing.
It says: to reconstruct $\xb$, compute its inner product with each basis
vector (the \emph{analysis} step), then build $\xb$ back up by combining
the basis vectors weighted by those inner products (the \emph{synthesis}
step).  In an orthonormal basis, analysis and synthesis are exact
inverses of one another.

\marginnote{%
When the basis is not orthonormal, equation~\eqref{eq:coord_orth} still
makes sense as the \emph{projection} of $\xb$ onto the direction of
$\mathbf{b}_k$, but the projections no longer give exact coordinates ---
the basis vectors interfere.  An orthonormal basis decouples them.
This is why so many useful representations (Fourier, the standard
basis, principal components) are designed to be orthonormal.}

We will see many examples of \eqref{eq:reconstruction} in disguise.
The Fourier transform is this formula with $\mathbf{b}_k$ a sinusoid of
frequency $k$.  The Discrete Wavelet Transform is this formula with
$\mathbf{b}_k$ a scaled, shifted copy of a mother wavelet.  Principal
component analysis is this formula with $\mathbf{b}_k$ the $k$-th
eigenvector of the data covariance.  A neural network with linear
activations is this formula with $\mathbf{b}_k$ a learned weight vector.

\emph{The choice of basis is the design of the representation.}  When we
pick sinusoids, we have decided that frequency is the property we want
to expose.  When we pick wavelets, we have decided that local
time-frequency structure matters.  When we let a network learn its
basis, we have decided to let the data decide.  Each choice has
consequences: it determines what the representation makes obvious, what
it hides, and what kinds of operations are easy versus hard to express
in it.  Part~\textsc{II} of this book builds the tools to make such
choices; Part~\textsc{III} explores their compositional consequences.

%------------------------------------------------------------------------
\section{Analysis and synthesis as matrices}
%------------------------------------------------------------------------

Equations~\eqref{eq:coord_orth} and~\eqref{eq:reconstruction} both
become cleaner when written in matrix form.  Let $\mathbf{B}$ be the
matrix whose $k$-th \emph{row} is the basis vector $\mathbf{b}_k^\top$.
Then computing all the inner products $\inner{\xb}{\mathbf{b}_k}$ at
once is the same as multiplying $\xb$ by $\mathbf{B}$:
\begin{equation}
    \mathbf{c} \;=\; \mathbf{B}\, \xb \qquad \text{(analysis)}.
    \label{eq:analysis_matrix}
\end{equation}
The vector $\mathbf{c}$ contains all the coordinates $c_k$.  Equivalently
written, $\mathbf{c} = \Phi(\xb)$ where $\Phi$ is the \emph{analysis
operator}.

For the synthesis direction, if the basis is orthonormal, the inverse of
$\mathbf{B}$ is its transpose (or, in the complex case, its Hermitian
adjoint $\mathbf{B}^H$).  Multiplying~\eqref{eq:analysis_matrix} on the
left by $\mathbf{B}^H$ recovers $\xb$:
\begin{equation}
    \xb \;=\; \mathbf{B}^H \mathbf{c} \;=\; \mathbf{B}^H\, \mathbf{B}\, \xb
    \qquad \text{(synthesis)}.
    \label{eq:synthesis_matrix}
\end{equation}
For non-orthonormal but still invertible $\mathbf{B}$, the synthesis
uses $\mathbf{B}^{-1}$ rather than $\mathbf{B}^H$.

So the entire analysis-synthesis cycle is just a pair of matrix-vector
multiplications.  This is the bridge to the rest of the book.  The
Fourier transform is exactly this with one specific choice of
$\mathbf{B}$ (a matrix of complex sinusoids).  The DWT is this with a
different specific choice (a matrix of wavelets at multiple scales).  A
linear neural-network layer is this with $\mathbf{B}$ \emph{learned from
data} instead of designed in advance.

%------------------------------------------------------------------------
\section{From basis to representation}
\label{sec:representations}
%------------------------------------------------------------------------

The basis idea generalises in three important ways, and the
generalisations are what make modern representations work.

\paragraph{Relaxing linear independence.}
In a classical basis, we insist that no basis vector is a linear
combination of the others.  This is what makes the representation
\emph{unique}: every $\xb$ has exactly one set of coordinates.  But
uniqueness is not always what we want.  Sometimes we want a redundant
set of vectors --- more than $N$ of them in an $N$-dimensional space ---
because the redundancy can encode information about the signal that a
minimal basis cannot.  Such sets are called \emph{frames} or
\emph{dictionaries}.  Matching pursuit, which we will meet in
Part~\textsc{III}, is a method for choosing a sparse subset of a
dictionary to represent a signal.

\paragraph{Allowing dimensionality change.}
In a classical basis, the matrix $\mathbf{B}$ is square: $N$ basis
vectors in an $N$-dimensional space.  If we allow $\mathbf{B}$ to be
rectangular --- say $K \times N$ with $K < N$ --- we get a
\emph{projection onto a subspace}.  The output has fewer dimensions than
the input.  This is dimensionality reduction, and it is what PCA, the
encoder of an autoencoder, and the early layers of a classifier all do.
The inverse direction --- $N$ to $K$ with $K > N$ --- is dimensionality
expansion, which is what the decoder of an autoencoder does, and what
generative models do when they turn a small latent vector into a full
image or audio file.

\paragraph{Nonlinearity.}
A pure projection $\xb \mapsto \mathbf{B}\xb$ is a linear operation;
composing two linear operations gives another linear operation, so
\emph{stacking} pure projections gives nothing new --- any depth of
linear layers collapses to a single equivalent layer.  To get a deep
representation that is more expressive than a shallow one, we need to
insert a nonlinearity between the projections.  In modern neural
networks, this is the ReLU function $\max(0, \cdot)$, the sigmoid
$\sigma$, or any of a small zoo of others.  In the Scattering Transform
(Chapter~\ref{ch:apriori}), it is the absolute value $\abs{\cdot}$ that
plays this role.

Putting all three together --- redundant dictionaries, varying
dimensionality, nonlinear stages --- gives the basic shape of a modern
representation:
\begin{equation}
    \mathbf{z}^{(k)} \;=\; f_k\!\left(\mathbf{B}_k\, \mathbf{z}^{(k-1)} + \mathbf{b}_k\right),
    \qquad k = 1, 2, \ldots, L,
    \label{eq:layered}
\end{equation}
with $\mathbf{z}^{(0)} = \xb$ the input, $\mathbf{B}_k$ the projection
matrix of the $k$-th layer, $\mathbf{b}_k$ a bias vector (allowing the
layer to be \emph{affine} rather than purely linear, which lets the
representation work with non-zero-centred data), and $f_k$ the
nonlinearity.

\marginnote{%
\textbf{Encoders and decoders.}  When the layers reduce dimensionality
($\mathbf{B}_k$ progressively narrower) and end in a small latent
vector, we call the whole stack an \emph{encoder}.  Reverse it ---
layers that progressively expand dimensionality from a latent vector
back to a full signal --- and we have a \emph{decoder}.  An
autoencoder is an encoder followed by a decoder, trained so that the
output reconstructs the input.}

Equation~\eqref{eq:layered} is the structural backbone of essentially
every neural network in this book.  Chapter~\ref{ch:learned} discusses
how the $\mathbf{B}_k$ and $\mathbf{b}_k$ are \emph{learned} from data
by gradient descent on a loss function.  In Part~\textsc{II} we will
build the code that does this.

%------------------------------------------------------------------------
\section{What we have so far, and what comes next}
%------------------------------------------------------------------------

We have set up the conceptual frame:

\begin{itemize}
    \item Signals are vectors in a finite-dimensional Hilbert space.
    \item The inner product measures global similarity between vectors.
    \item Convolution is the inner product applied locally, kernel-by-kernel.
    \item Representations are projections onto a chosen basis (or
    dictionary), and reconstruction is the reverse operation.
    \item The choice of basis is the design of the representation; it
    determines what is made obvious and what is hidden.
    \item Layered representations stack projections separated by
    nonlinearities, allowing them to express far more than a single linear
    projection can.
\end{itemize}

\noindent The next two chapters give specific instances.
Chapter~\ref{ch:apriori} treats \emph{a-priori} representations --- those
whose basis is designed in advance --- with the Fourier transform, the
wavelet transform, and the Scattering Transform as our three main
examples.  Chapter~\ref{ch:learned} treats \emph{learned} representations
--- those whose basis is derived from data --- with principal component
analysis, $k$-means clustering, matching pursuit, and neural networks.
Chapter~\ref{ch:style_transfer} closes Part~\textsc{I} by working a
single concrete problem --- musical style transfer --- through each of
these representations in turn, showing how the choice changes what is
possible.  This sets up the technical work of Part~\textsc{II}, where we
build the code, and the compositional work of Part~\textsc{III}, where
we use it.
